\begin{proposition}[Proved: a large leaf in a small three-object star]
\label{CP2:A-star-large-leaf}
Let $C$ be a connected three-object category whose acyclic support
is a star, with no chain of three distinct objects. If $C$ has at
most seventeen arrows, identities included, and a leaf corner has
order at least eight, then $BC$ is not homotopy equivalent to
$\mathbb{CP}^2$.
\end{proposition}
\begin{proof}
After taking the opposite category, write the two nonempty cross-homs
as $0\to1$ and $0\to2$. Let the corner orders be $a,b,c$, with
$b\geq8$, and let the stalk sizes be $s,t$. Thus
\[
 a+b+c+s+t\leq17,\qquad a,b,c,s,t\geq1.
\]
The center order satisfies $a\leq6$. Let $E_1,E_2$ be the full
edge subcategories on $\{0,1\}$ and $\{0,2\}$. Every nerve simplex
lies in one edge subcategory, so $BC=BE_1\cup_{BM}BE_2$, where
$M=C(0,0)$. These are CW subcomplexes and this union is their
homotopy pushout.

If $a=6$, the budget forces $b=8$ and $c=s=t=1$. All cross-homs
are singletons, so Lemma~\ref{CP2:A-singleton-support} identifies
rational second cohomology with that of the support tree. It is
zero. If $a=5$, both $s,t$ are at most two. A singleton cross-hom
is the join of the two corner classifying spaces, and has zero
rational second cohomology: both factors are connected and finite
monoids have zero rational first homology. A two-arrow cross-hom
has zero rational degree-two products by
Lemma~\ref{CP2:A-high-two-cross}. Also $H^3(BM;\mathbb Q)=0$ by
Corollary~\ref{CP2:A-five-types}. Mayer--Vietoris therefore makes
\[
 H^4(BC;\mathbb Q)\longrightarrow
 H^4(BE_1;\mathbb Q)\oplus H^4(BE_2;\mathbb Q)
\]
injective. Every global degree-two product restricts to zero in
both components, hence is zero.

It remains that $a\leq4$. Choose an idempotent $e$ in the minimal
ideal $I$ of $M$ and write $G=eMe$. The small-minimal-ideal proof
in Lemma~\ref{CP2:A-small-minimal-ideal} shows that at least one
of $Me=G$ and $eM=G$ holds, and that $BM\simeq BG$.
If $G$ is trivial, $BM$ is contractible. Collapsing this CW
subcomplex gives a homotopy equivalence
\[
 BC\simeq BE_1\vee BE_2.
\]
Indeed the common contractible subcomplex is collapsed in each
edge as well, and those quotient maps are homotopy equivalences.
If the wedge had type $\mathbb{CP}^2$, its free-product fundamental
group would force both edge spaces to be simply connected. Its
integral homology is the direct sum of theirs in positive degrees.
The nonzero primitive square forces the rank-one second and fourth
homology to lie in the same edge summand. That edge has exactly
the integral homology and primitive cup square of $\mathbb{CP}^2$,
and no other reduced homology. The recognition theorem makes it
an actual projective-plane model. Its arrow count is at most
$17-1-1=15$, contradicting the newly proved two-object baseline,
Theorem~\ref{CP2:A-two-object-baseline-sixteen}. This argument uses
that theorem for an actual homotopy type, not merely a rational
cohomology pattern.

Suppose finally that $G$ is nontrivial. In this small order both
$Me=G$ and $eM=G$ must hold. For example, if $a_0\in Me\setminus G$,
the map $G\to a_0G$, $g\mapsto a_0g$, is injective: multiplication
by $e$ reduces equality to equality in $G$, since $ea_0\in G$.
The sets $G$ and $a_0G$ are disjoint and lie in $I$, while the global
identity lies outside $I$ unless $M$ itself is a group. This would
require at least $2|G|+1\geq5$ elements. The other side is dual.
It follows that $I=MeM=G$. The averaging idempotent
\[
 f=|G|^{-1}\sum_{g\in G}g\in\mathbb QM
\]
satisfies $mf=f=fm$ for every $m\in M$: multiplication by $m$
acts on the ideal group by a group translation. Thus
$(\mathbb QM)f$ and $f(\mathbb QM)$ are the one-dimensional trivial
left and right modules, respectively, and both are projective.

Collapse the center to a trivial object, and replace each incident
stalk by its actual $M$-orbit set. Precomposition gives each stalk
its right $M$-action; the commuting leaf action descends to the
orbit set. These assignments define an actual quotient functor.
Lemma~\ref{CP2:A-high-orbits} makes it a rational homology and
cohomology-ring equivalence in all degrees. Its target $C'$ has
trivial center, hence
\[
 BC'=BX_1\vee BX_2,
\]
where $X_1,X_2$ are ordinary act cones, up to opposition, over the
leaf monoids. Their state counts are at most $s,t$. If $BC$ had
type $\mathbb{CP}^2$, this wedge would have its rational homology
and cohomology ring, and both summands would have bounded rational
homology.

Here $a\geq2$, so the smaller leaf has
$c\leq17-a-b-2\leq5$. Its act cone has zero rational degree-two
products by Theorem~\ref{CP2:A-bounded-nine-acts}. For the large
leaf, the number of states is at most
\[
 s\leq17-a-b-c-t\leq13-b.
\]
For $b=8$ or $9$, the same bounded-nine theorem applies. For
$b\geq11$, there are at most two states, so
Lemma~\ref{CP2:A-high-affine} gives zero products. The only remaining
case is $b=10$, with at most three states. Since the other wedge
summand has zero products, the nonzero square of the global
rank-one degree-two class would lie in $BX_1$. The direct-sum
homology then gives $BX_1$ exactly the rational homology and
cohomology ring of $\mathbb{CP}^2$. This contradicts
Proposition~\ref{CP2:A-ten-six-state-rational}. Every case is excluded.
\end{proof}
